\documentclass[9pt,letterpaper]{article}
\usepackage[utf8]{inputenc}
\usepackage[T1]{fontenc}
\usepackage{geometry}
\geometry{top=0.55in,bottom=0.55in,left=0.68in,right=0.68in}
\usepackage{enumitem}
\setlist{itemsep=0pt, topsep=1pt, parsep=0pt}
\usepackage{hyperref}
\usepackage{setspace}
\usepackage{siunitx}
\sisetup{separate-uncertainty=true, range-phrase=--}
\DeclareSIUnit\angstrom{\text{\AA}}

\begin{document}
\fontsize{9}{10.5}\selectfont

\begin{flushleft}
\textbf{Myke Vital Sao Temgoua} \\
Department of Physics, Faculty of Science \\
University of Yaound\'{e} I \\
P.O. Box 812, Yaound\'{e}, Cameroon \\
\href{mailto:myke-vital.sao@facsciences-uy1.cm}{myke-vital.sao@facsciences-uy1.cm} \\[0.8em]
\today
\end{flushleft}

\begin{flushleft}
Prof.\ Kenneth M.\ Merz, Jr., Editor-in-Chief \\
\textit{Journal of Chemical Information and Modeling} \\
American Chemical Society
\end{flushleft}

\vspace{0.2em}

\noindent Dear Prof.\ Merz,

We submit our manuscript \textbf{``Target breadth and mutation resilience in African-natural-product-inspired antimalarial chemotypes: a computational analysis''} for consideration as an Article in \textit{Journal of Chemical Information and Modeling}.

This work addresses a critical methodological challenge in computational antimalarial drug discovery: evaluating target breadth and mutation resilience as distinct, complementary properties when prioritizing candidates from African natural-product-inspired chemical space, while maintaining clear boundaries between computational evidence and biological activity. We analyze a 17-member polypharmacology-oriented cohort across four mechanistically diverse antimalarial targets (PfDHFR, PfCRT, PfClpP, and PfATP4) using target-specific structural anchors and a per-target resistance-resilience scoring (RRS) framework. The workflow integrates three evidence layers: chemical-space expansion (\num{65856} hybrid library molecules $\rightarrow$ \num{19913} synthesizable candidates), target-anchored AutoDock Vina docking generating a complete $\num{17}\times\num{4}$ scoring matrix (\num{68}/\num{68} geometric quality pass rate), and per-target RRS calculated separately for PfDHFR (N51I, C59R, S108N, I164L) and PfCRT (K76T, K76A) mutant panels, with PfCRT re-docked against a corrected 3D7-like wild-type receptor under a cavity-anchored grid and a pipeline-null control demonstrating that the corrected PfCRT channel carries no detectable K76T/K76A signal for this panel. The cohort yielded seven A*, four B, and six C RRS profiles (range \numrange{73.2}{115.6}); after multiplicity correction, PNS--RRS ($\rho=-\num{0.714}$), RRS--wild-type score ($\rho=\num{+0.691}$), and $N_{\mathrm{fav}}$--RRS ($\rho=\num{0.714}$) were significant, while ACSI--RRS was not ($\rho=-\num{0.190}$).

\noindent\textbf{Key enhancements} (comprehensive validation, physicochemical characterization, and methodological documentation):

\begin{enumerate}
\item \textbf{Validation documentation:} three complementary approaches: (a) external benchmark (DEKOIS 2.0, PfDHFR ROC-AUC \num{0.45}, reported honestly as near-chance), (b) enrichment validation (MMV Malaria Box, ROC-AUC \numrange{0.924}{1.000}), and (c) redocking validation (four of five reference ligands reproduced with RMSD $< \SI{2.0}{\angstrom}$; the methotrexate exception reflects NADPH-adjacent grid placement). Metrics in Supporting Information Section S12.

\item \textbf{Retrospective boundary check:} RRS profiles computed for five approved antimalarials (artemisinin, pyrimethamine, chloroquine, lumefantrine, cipargamin) against the same PfDHFR and corrected PfCRT panels: the framework does not recover the clinically documented resistance signatures (pyrimethamine $\approx$ \qty{100}{\percent} on all four PfDHFR mutants; chloroquine $\approx$ \qty{99}{\percent} on PfCRT K76T/K76A, both within the pipeline-null spread), reinforcing that docking-RRS values are not resistance-circumvention evidence (Supporting Information Table S7).

\item \textbf{Physicochemical characterization:} all 17 candidates pass the reported Lipinski/Veber criteria (\qty{100}{\percent} compliance; mean QED \num{0.703}); the available ADMET summary is population-level and is not a candidate-specific safety assessment. Details in Supporting Information Tables S2--S4.

\item \textbf{Multi-parameter optimization framework:} a five-component MPO (Vina affinity 35\%, DiffDock confidence 25\%, QED 20\%, ADMET 15\%, Lipinski penalty 5\%) with a polypharmacology bonus and an MPO $\geq$ 0.40 selection threshold (Section S10).

\item \textbf{Methodological rigor:} complete grid box specifications for all four targets, \qty{99.3}{\percent} cost reduction via centroid-based library reduction, and strict evidence boundaries.
\end{enumerate}

\noindent\textbf{Evidence boundaries:} the study does not claim experimental affinity, confirmed simultaneous target engagement, biological resistance circumvention, or measured IC\textsubscript{50}/EC\textsubscript{50} values. The workflow produces (a) four-target docking profiles from which polypharmacology hypotheses can be generated, (b) per-target RRS values summarizing mutation-panel docking ratios, and (c) exploratory cross-metric associations without conflation.

\noindent\textbf{Reproducibility and impact:} complete computational provenance (source code, candidate manifests, receptor/ligand files, grid configurations, raw poses, RRS protocols, machine-readable tables) is archived at \url{https://github.com/NanaEngo/Malaria_codesV2}; the Supporting Information (\num{22} pages, \num{13} sections, \num{14} tables, \num{2} figures) contains the validation datasets and methodological specifications.

The centroid-based reduction strategy lowers the computational barrier for endemic-region groups (\qty{94}{\percent} of malaria cases) by orders of magnitude while retaining structural diversity, providing a reproducible template for hypothesis generation in resistance-challenged therapeutic areas. This manuscript is original, not under consideration elsewhere, and all authors have reviewed and approved the submission. We declare no competing financial interests.

\vspace{0.3em}

\noindent Sincerely,

\noindent \textbf{Myke Vital Sao Temgoua} (on behalf of all co-authors)

\end{document}
